\scanpage{132}

previously concerning the $N_V,I_V\ldots.$), I am going to look a little more closely at the form of the subgroups $I_x\subset{\unsure{\pi|\mathcal E}}_{L/K}^{\bar L}$, or \deleted{\unsure{associated}} rather, of the \inserted{$\pi$-}conjugacy classes of such subgroups) associated with the $x\in X\setminus U$, \inserted{for varying $x$,} and convince myself of the intrinsic character of these groups, independently of the choice of the model $X$ (but this is obvious, if one admits what one has just \unsure{stated|announced} --- that they are precisely, up to open subgroups \deleted{\unsure{all|the}} at most, the images of the $I_V$, for $V\deleted{\unsure{\to}}\inserted{\in}\mathfrak X_{L/K}$ not having \inserted{a} center in $\deleted{X\setminus U}U\ldots$)

\deleted{Let us first consider the}

Let $J$ be a subset of the finite set $I$, and let
\[
\tag{36}
D_J^*=\{x\in X\mid I(x)=J\}
=\overbrace{\bigcap_{j\in J}D_j}^{\Delta_J}-\bigcup_{i\in I\unsure{\setminus|-}J}\Delta_J\cap D_i
\]
so that
\[
D_{\{i\}}=D_i\setminus\left(\bigcup_{j\ne i}\unsure{D_i\cap D|D_i\cap D_j}\right).
\]
If $d=\operatorname{card}J$, \unsure{one} sees that $D_J\subset X$ is \inserted{regular,} of codimension\textsuperscript{\unsure{n}} $d$ at each of its points (but of course, it may be $\phi$). If one calls a subset $J\subset I$ a ``simplex\unsure{''|'}

\margin{it is better to define $D_J$ as a \uline{\unsure{subscheme|subscheme}} as $\subset\bigcap_{j\in J}D_j$, and define its \uline{support} as $\{x\in X\mid I(x)=J\}$}

\margin{131}
